% TOP_PROBLEM: {"id": 30006694, "title": "The Exponential Time Hypothesis for 3-SAT", "category_id": 15, "category": {"id": 15, "name": "computer_science", "display_name": "Computer Science", "description": "Computational complexity, algorithms, and theoretical CS.", "slug": "computer-science", "order_index": 15, "created_at": "2026-07-31T15:26:25.671Z"}, "status": "open"}
% FIRST_POSTED: 2026-09-27
% !TeX program = lualatex
% Compile twice: lualatex open_problems_500.tex
% All references and prose are embedded; no BibTeX database or external inputs.
\documentclass[11pt,a4paper]{article}
\usepackage[margin=24mm,headheight=14pt]{geometry}
\usepackage{fontspec}
\setmainfont{Latin Modern Roman}
\setsansfont{Latin Modern Sans}
\setmonofont{Latin Modern Mono}
\usepackage{amsmath,amssymb,mathtools}
\usepackage{unicode-math}
\setmathfont{Latin Modern Math}
\usepackage{microtype}
\usepackage{enumitem,xurl,needspace}
\usepackage[dvipsnames]{xcolor}
\usepackage{hyperref}
\hypersetup{unicode=true,colorlinks=true,linkcolor=MidnightBlue,urlcolor=MidnightBlue,
 pdftitle={Research notebook: Problems 40--49},pdfauthor={Source-annotated compilation},bookmarksnumbered=true}
\usepackage{bookmark}
\usepackage{tocloft}
\setlength{\cftsecnumwidth}{3em}
\cftsetpnumwidth{2.5em}
\cftsetrmarg{4em}
\usepackage{titlesec}
\titleformat{\section}{\Large\bfseries\raggedright}{\thesection}{0.7em}{}
\usepackage{fancyhdr}
\pagestyle{fancy}\fancyhf{}
\fancyhead[L]{\small\sffamily Mathematical problem statements}
\fancyhead[R]{\small Source edition 22}
\fancyfoot[C]{\thepage}
\setlength{\parindent}{0pt}
\setlength{\parskip}{5pt plus 1pt}
\setlength{\emergencystretch}{3em}
\setcounter{tocdepth}{1}
\setcounter{secnumdepth}{1}
\setlist[itemize]{leftmargin=3em,itemsep=5pt,topsep=4pt,parsep=0pt}
\allowdisplaybreaks
\newcommand{\N}{\mathbb{N}}
\newcommand{\Z}{\mathbb{Z}}
\newcommand{\Q}{\mathbb{Q}}
\newcommand{\R}{\mathbb{R}}
\newcommand{\C}{\mathbb{C}}
\newcommand{\F}{\mathbb{F}}
\newcommand{\A}{\mathbb{A}}
\newcommand{\PP}{\mathbb P}
\newcommand{\E}{\mathbb{E}}
\newcommand{\eps}{\varepsilon}
\newcommand{\dd}{\,\mathrm d}
\newcommand{\sref}[2]{\hyperlink{source:#1:#2}{[S#2]}}
\newcommand{\eref}[2]{\hyperlink{extra:#1:#2}{[E#2]}}
% Turn the next switch off for a shorter reading copy. The quotations preserve
% every exactTarget field, including qualifications omitted from a short summary.
\newif\ifshowcatalogue
\showcataloguetrue
\newcommand{\cataloguescope}[1]{\ifshowcatalogue
 \paragraph{Exact scope in the supplied catalogue.}
 \begin{quote}\small #1\end{quote}\fi}

% Research additions dated 27 September 2026; compile twice with LuaLaTeX.
\numberwithin{equation}{section}
\begin{document}
\setcounter{section}{0}
% ================================================================
% problemId: 30006694
% Canonical problem ID: 30006694
% exactTarget SHA-256: 31392dee0d6193401d7cab8b03526893ec5658995559efc42735f2f719d4d30a
\clearpage
\section*{\texorpdfstring{The Exponential Time Hypothesis for 3-SAT}{The Exponential Time Hypothesis for 3-SAT}}\label{30006694}
\subsection{Definitions and mathematical statement}
A $3$-CNF formula is a conjunction of clauses, each a disjunction of at most three literals in Boolean variables. Let $n$ be the number of variables and $|F|$ the input length. ETH asserts a strictly positive exponential-time rate for deciding satisfiability, conventionally expressed as
\[
 \exists c>0:\quad 3\text{-SAT has no deterministic }O(2^{cn}\operatorname{poly}(|F|))\text{-time algorithm}.
\]
Polynomial input-reading factors and the precise time-bound convention are implicit in the catalogue's abbreviated $2^{cn}$ wording. The parameter is the number of variables, not the numerical size of an encoding or the clause count alone.
\par\smallskip\textit{Formulation and concept references:} \sref{30006694}{1}.
\cataloguescope{Prove or disprove the Exponential Time Hypothesis: there exists a constant c > 0 such that 3-SAT cannot be solved in deterministic time 2\textasciicircum{}\{c n\}.}
\subsection{Short English statement}
Does some positive exponential dependence on the number of variables remain unavoidable for deterministic $3$-SAT algorithms?
% BEGIN RESEARCH ADDITION 45
\subsection{Research attempt}
\paragraph{Current frontier.}
The positive lower bound on the exponential rate in \sref{30006694}{1} is
stronger than $\mathsf P\ne\mathsf{NP}$, as also emphasized by the
2026 notes \sref{30006694}{2}. The sparsification lemma of
Impagliazzo--Paturi--Zane reduces general fixed-width SAT to a controlled
number of bounded-occurrence instances \eref{30006694}{1}. A published
benchmark deterministic algorithm for $3$-SAT has running time
$O^*(1.32793^n)$ \eref{30006694}{2}; quoting such an upper bound does not prove
that every algorithm needs any particular positive exponent. Recent
strong sparsification by merging variables concerns $1$-in-$3$-SAT,
with its different acceptance predicate, and cannot simply be substituted
for ordinary $3$-SAT \eref{30006694}{4}. The attempt below therefore looks
for a subexponential algorithm, then tests the structural lemmas that
would make it work.

\paragraph{Attack route.}
The first candidate is a sparse-instance separator algorithm: enumerate
assignments on a small separator and solve independent residual pieces.
The precise missing property is an efficiently discoverable, recursively
sublinear boundary after sound preprocessing. A second candidate replaces
large boundary tables by compressed functions. That requires an exact
representation closed under conjunction and existential quantification,
with subexponential representation size and update time. A third route,
proving ETH from large tables or large matrix ranks, would need a lower
bound applying to every deterministic computation, not merely to a
particular decomposition. We investigate the separator route and
explicitly invalidate the simplest rank-based converse.

\paragraph{Attempt: a separator-to-algorithm lemma.}
The primal graph of a CNF formula has one vertex per variable and an edge
between variables occurring in a common clause. Consider a hereditary
class of such graphs for which a deterministic polynomial-time procedure
finds, on every $v$-vertex induced subgraph, a set $S$ of size at most
$s(v)$ such that all components after deletion have at most $2v/3$
vertices. Suppose $s(v)=o(v)$, uniformly over that class.

Recursively delete separators until only bounded-size pieces remain.
For each node make a bag containing its separator and all ancestor
separators; leaves also contain their remaining variables. This is a
tree decomposition: every vertex has connected bag occurrences, and an
edge cannot connect two different components below the separator that
first separates them. Every clause is contained in some bag, since its
variables form a clique. The largest bag size $w+1$ is bounded by the
sum of separator sizes along one root-to-leaf path, plus the leaf size.

Here the little-oh estimate really is uniform. Given $\epsilon>0$, choose
$V_\epsilon$ so that $s(v)\le\epsilon v$ for $v\ge V_\epsilon$.
Along a path the subproblem sizes decrease by a factor at most $2/3$.
The contribution before reaching $V_\epsilon$ is at most
\[
 \epsilon n\sum_{j\ge0}(2/3)^j=3\epsilon n.
\]
The remaining path has bounded length and bounded bag contribution,
depending on $V_\epsilon$ but not on $n$. Hence $w=o(n)$.

For completeness, SAT can be decided from this decomposition without a
probabilistic step. At each bag store which Boolean assignments to the
bag extend to satisfying assignments of all clauses allocated to its
subtree. Check the clauses allocated to the bag and combine child tables
by agreement on overlaps and existentially forgetting child-only
variables. The separator property ensures there are no missing
cross-child clauses. There are $2^{w+1}$ assignments per bag; even a
naive pairwise implementation costs only $2^{O(w)}$ times a polynomial
in the input and decomposition size. Thus the stated separator hypothesis
gives a deterministic $2^{o(n)}\operatorname{poly}(|F|)$ algorithm on the
class. This is an auxiliary algorithmic lemma, not a proof that all
sparse formulas satisfy its hypothesis.

\paragraph{Attempt: the exact rate calculation with sparsification.}
The sparsification theorem \eref{30006694}{1}, Section 2 and Corollary 1,
says that for each fixed $\alpha>0$ a $3$-CNF on $n$ variables can be
expressed as the disjunction of at most $2^{\alpha n}$ formulas on those
variables, each with occurrence bound depending only on $\alpha$.
The list is generated in time $2^{\alpha n}\operatorname{poly}(|F|)$.
In particular every output formula has at most $D_\alpha n$ clauses,
for a constant $D_\alpha$.

Suppose one could solve every fixed-density class $m\le Dn$ in time
$2^{o_D(n)}\operatorname{poly}(|F|)$, for example by a valid extension
of the preceding separator method. Solving all sparsified formulas would
cost
\[
 2^{\alpha n+o_{D_\alpha}(n)}\operatorname{poly}(|F|). \tag{45.1}
\]
For any proposed ETH constant $c>0$, fix $\alpha=c/3$. For sufficiently
large $n$, the little-oh term is at most $cn/3$, so (45.1) is bounded
by $2^{2cn/3}\operatorname{poly}(|F|)$, with the finitely many smaller
inputs absorbed into a constant. This supplies an algorithm within the
rate excluded by that $c$. Consequently such a collection of sparse
algorithms would refute the \emph{exact existential-rate formulation}
of ETH in the original entry.

The quantifiers here are important: the selected algorithm can depend
on the fixed $c$, not on the growing input. We have not silently asserted
that a non-effectively chosen family of algorithms automatically becomes
one algorithm with a uniformly specified $2^{o(n)}$ running time. The
rate calculation already contradicts the displayed target without that
extra uniformization claim.

\paragraph{Stress test: sparse graphs need not have small separators.}
Bounded-degree vertex-expander families exist; see \eref{30006694}{3}, Sections
2 and 4. Let $G$ be such a graph with maximum degree $d$ and expansion
constant $h>0$, meaning every $U$ with $|U|\le n/2$ has at least
$h|U|$ neighbors outside $U$. A separator $S$ whose remaining components
have size at most $2n/3$ must have linear size. If $|S|\ge n/6$ this is
immediate. Otherwise a union of components can be chosen with size
between $n/6$ and $n/2$: choose one in that range, greedily combine
smaller ones, or take the complement of a component larger than $n/2$
inside $G-S$. All external neighbors of this union lie in $S$.
Therefore
\[
 |S|\ge hn/6\quad\hbox{or}\quad |S|\ge n/6. \tag{45.2}
\]
This disproves the universal graph-theoretic separator hypothesis even
when the clause count is linear. Indeed
\[
 F_G=\bigwedge_{\{u,v\}\in E(G)}(x_u\lor x_v) \tag{45.3}
\]
is a $3$-CNF in the entry's ``at most three literals'' convention, has
at most $dn/2$ clauses, and has primal graph $G$. Yet (45.3) is
satisfied by the all-one assignment and can be recognized as such
immediately. Thus large separators do not prove computational hardness
either. A successful extension would have to exploit the actual
predicates or a sound semantic simplification, not just the primal graph.

A simple preprocessing calculation illustrates the remaining tradeoff.
For a formula with $m\le Dn$, let $T$ contain variables with more than
$\Delta$ occurrences. Since total occurrences are at most $3Dn$,
$|T|\le3Dn/\Delta$. Enumerating $T$ costs at most
$2^{3Dn/\Delta}$ branches, after which every remaining variable has at
most $\Delta$ occurrences. This reduces high occurrence but leaves the
bounded-degree obstruction (45.2); it does not produce sublinear
separators. It cannot be promoted to a subexponential algorithm by
ignoring the cost of solving the residual instances.

\paragraph{Stress test: an exponential boundary rank is not ETH.}
Consider $2b$ variables and the formula
\[
 E_b(x,y)=\bigwedge_{i=1}^b
    [(\neg x_i\lor y_i)\land(x_i\lor\neg y_i)]. \tag{45.4}
\]
The truth table across the cut $x\mid y$ is the $2^b$ by $2^b$
identity matrix, since (45.4) holds precisely when $x=y$. It has rank
$2^b$ over every field. Any representation as a sum of separated scalar
products $\sum_{j=1}^r a_j(x)b_j(y)$ therefore requires $r\ge2^b$.
Nevertheless the formula is trivially satisfiable, describes only
pairwise equalities, and has a width-one decomposition if variables are
ordered as pairs. Its matrix-free description has linear size.

The exact ranks for $b\le6$ are checked in the supplied script. The
calculation refutes a lower-bound argument based on one chosen cut; it
is not evidence against a better optimized representation. To prove ETH
one would need to rule out all deterministic algorithms, including
algorithms that never materialize that table. To disprove ETH using
compression, one must give a uniform update procedure and account for
all existential projections; naming a short formula for the original
input does not show that its projections stay short.

\paragraph{Outcome.}
\textbf{Outcome: Exploratory approach with unresolved gap.}
The attempt proves a separator-based subexponential algorithm under a
precise hereditary hypothesis, and shows exactly how sparse-instance
algorithms with vanishing rate would contradict the stated ETH
quantifiers. The unqualified separator hypothesis is false, as shown
by explicit easy formulas on expanders. A separate exact construction
shows why exponential boundary rank does not furnish an unrestricted
time lower bound. These are auxiliary deductions, not assertions of
historical novelty or a resolution of ETH. The unresolved step is a
semantic decomposition or compression theorem handling every sparse
formula with a genuinely vanishing exponential rate; no such theorem
is supplied. Conversely no universal positive-rate lower bound is
established. A proof of ETH would in particular prove $\mathsf P\ne\mathsf{NP}$.
% END RESEARCH ADDITION 45
\subsection{Sources}
\begin{itemize}\raggedright
\item[\textnormal{[S1]}]\hypertarget{source:30006694:1}{} R. Impagliazzo, R. Paturi, J. Comput. System Sci. 62(2):367-375, 2001.
\url{https://doi.org/10.1006/jcss.2000.1727}.
{\small\textit{Catalogue formulation source.}}
\item[\textnormal{[S2]}]\hypertarget{source:30006694:2}{} Complexity Theory: Stronger forms of P != NP, Roei Tell.
\url{https://www.cs.toronto.edu/~roei/2026%20complexity/6%20stronger%20lbs.pdf}.
{\small\textit{Catalogue research/status reference; not a proof endorsement.}}
% BEGIN RESEARCH SOURCES 45
\item[\textnormal{[E1]}]\hypertarget{extra:30006694:1}{} Russell Impagliazzo, Ramamohan Paturi and Francis Zane,
\emph{Which Problems Have Strongly Exponential Complexity?}, Journal of Computer and System Sciences
63 (2001), 512--530, Section 2, especially Corollary 1 (sparsification).
\url{https://cseweb.ucsd.edu/~russell/ipz.pdf}.
{\small\textit{Primary sparsification theorem and its uniform parameter dependence. Consulted 27 September 2026.}}
\item[\textnormal{[E2]}]\hypertarget{extra:30006694:2}{} Sixue Liu,
\emph{Chain, Generalization of Covering Code, and Deterministic Algorithm for k-SAT},
ICALP 2018, LIPIcs 107, 88:1--88:13, abstract and the deterministic $3$-SAT bound.
\url{https://drops.dagstuhl.de/entities/document/10.4230/LIPIcs.ICALP.2018.88}.
{\small\textit{A published algorithmic benchmark, not asserted to be a newly verified optimal bound. Consulted 27 September 2026.}}
\item[\textnormal{[E3]}]\hypertarget{extra:30006694:3}{} Shlomo Hoory, Nathan Linial and Avi Wigderson,
\emph{Expander graphs and their applications}, Bulletin of the American Mathematical Society
43 (2006), 439--561, Sections 2 and 4.
\url{https://www.cs.huji.ac.il/~nati/PAPERS/expander_survey.pdf}.
{\small\textit{Existence and expansion properties of bounded-degree expanders. Consulted 27 September 2026.}}
\item[\textnormal{[E4]}]\hypertarget{extra:30006694:4}{} Benjamin Bedert, Tamio-Vesa Nakajima, Karolina Okrasa and Stanislav \v{Z}ivn\'y,
\emph{Strong Sparsification for 1-in-3-SAT via Polynomial Freiman-Ruzsa},
arXiv:2507.17878, full version of a FOCS 2025 paper, abstract.
\url{https://arxiv.org/abs/2507.17878}.
{\small\textit{Adjacent sparsification result with a different constraint predicate. Consulted 27 September 2026.}}
% END RESEARCH SOURCES 45
\end{itemize}

\end{document}
